% !TeX root = ../main.tex

\section{Assignment \#1}

\begin{problem}
  Please analyze the symmetry of the following shapes by
  \begin{enumerate*}
    \item listing \textbf{all the symmetry elements};
    \item using the flow chart to determine which point groups these structures
    belong to.
  \end{enumerate*}
  \begin{figure}[htbp]
    \begin{minipage}{.68\linewidth}
      \centering
      \includegraphics[width = \linewidth]{image7.png}
    \end{minipage}
    \hfill
    \begin{minipage}{.28\linewidth}
        \centering
        \includegraphics[width = .8\linewidth]{HW1-1.png}\\\medskip
        \includegraphics[width = .8\linewidth]{HW1-2.png}
    \end{minipage}
  \end{figure}
  \begin{enumext}
    \item A \textbf{planar} `cluster' of a tri-water-molecule `bonded' with
    hydrogen bonds between oxygen and hydrogen atoms (see the figure below for
    detailed structural parameters). (\textbf{0.5 point})
    \item Hexagonal Prism (It has two identical parallel hexagonal bases-top
    and bottom faces, and six rectangular faces connecting corresponding edges of the hexagons). (\textbf{0.5 point})
  \end{enumext}
\end{problem}
\begin{solution}\leavevmode
  \begin{enumext}
    \item
    \begin{enumext}
      \item Symmetry elements.
      \begin{enumext}[label = \arabic*.]
        \item Identity ($E$).
        \item Principal axis ($C_3$): An axis through the center of the hexagon
        perpendicular to the plane. Rotating the cluster by $\qty{360}\degree/3$
        maps each water molecule onto the next one.
        \item Horizontal mirror plane ($\sigma_h$): Since the cluster is planar,
        the plane containing all the atoms acts as a mirror plane.
        \item Improper rotation axis ($S_3$): Resulting from performing a $C_3$
        rotation followed by a reflection in the $\sigma_h$ plane.
      \end{enumext}
      \item Mapping the symmetry elements to the flow chart.
      \begin{enumext}[label = \arabic*., columns = 2]
        \item Nonlinear molecule.
        \item Only have one $C_3$ axis.
        \item $C_n$ ($n = 3 > 1$) satisfies.
        \item No $C_2$ axis.
        \item Horizontal mirror plane $\sigma_h$ exists.
      \end{enumext}
      \begin{center}
        \menu[,]{%
          Linear (\textcolor{red}{N}),
          $\geq 2$ non-overlapping $C_{n>2}$ (\textcolor{red}{N}),
          $C_n$ ($n > 1$) (\textcolor{red}{Y}),
          $nC_2$ normal to highest $C_n$ (\textcolor{red}{N}),
          $\sigma_h$ (\textcolor{red}{Y}),
          $C_{3h}$%
      }
      \end{center}
    \end{enumext}
    \pagebreak
    \item
    \begin{enumext}
      \item Symmetry elements.
      \begin{enumext}[label = \arabic*.]
        \item Identity ($E$).
        \item Principal axis ($C_6$): A $6$-fold rotation axis passing through
        the centers of the top and bottom hexagonal faces,
        \emph{coincident with} $C_3$ and $C_2$ rotations.
        \item Perpendicular $C_2$ axes ($6 \times C_2$): Three pass through
        the centers of opposite vertical rectangular faces, and three pass
        through opposite vertical edges.
        \item One horizontal mirror plane and six vertical mirror planes:
        ($\sigma_h$, $3 \times \sigma_v$, $3 \times \sigma_d$).
        \item Inversion center ($i$).
        \item Improper rotation axes ($S_6$, $S_3$, $S_2$).
      \end{enumext}
      \item Mapping the symmetry elements to the flow chart.
      \begin{enumext*}[label = \arabic*., columns = 2]
        \item(2) Nonlinear molecule.
        \item(2) No $C_{n>2}$ axes perpendicular to the principal axis:
        Only have $C_2$ perpendicular to it.
        \item $C_n$ ($n = 6 > 1$) satisfies.
        \item $6C_2$ axes perpendicular to $C_6$ axis.
        \item Horizontal mirror plane $\sigma_h$ exists.
      \end{enumext*}
      \begin{center}
        \menu[,]{%
          Linear (\textcolor{red}{N}),
          $\geq 2$ non-overlapping $C_{n>2}$ (\textcolor{red}{N}),
          $C_n$ ($n > 1$) (\textcolor{red}{Y}),
          $nC_2$ normal to highest $C_n$ (\textcolor{red}{Y}),
          $\sigma_h$ (\textcolor{red}{Y}),
          $D_{6h}$%
      }
      \end{center}
    \end{enumext}
  \end{enumext}
\end{solution}

\begin{problem}
  Continue working on the decomposition of the reducible representation
  of \ce{BF3} into irreducible representations (irreps., as various
  molecular modes).
  \begin{enumext}
    \item Point out which irreps. correspond to
    \begin{enumerate*}
      \item translational,
      \item rotational, and
      \item vibrational
    \end{enumerate*}
    modes. (\textbf{0.5 point})
    \begin{remark}[The \ce{BF3} ($D_3h$) point group character table]%
    \leavevmode\par
      \begin{center}
        \begin{tabular}{*9l}
          \toprule
          $D_{3h}$ & $E$         & $2C_3(z)$ & $3C_2'$ & $\sigma_h(xy)$      &
          $2S_3$   & $3\sigma_v$ & \parbox{7em}{Linear functions, rotations} &
          \parbox{4em}{Quadratic functions}\\
          \midrule
          $A_1'$   & $+1$        & $+1$      & $+1$    & $+1$                &
          $+1$     & $+1$        & ---       & $x^2 + y^2$, $z^2$            \\
          $A_2'$   & $+1$        & $+1$      & $-1$    & $+1$                &
          $+1$     & $-1$        & $R_z$     & ---                           \\
          $E'$     & $+2$        & $-1$      & $0$     & $+2$                &
          $-1$     & $0$         & $(x, y)$  & $x^2 - y^2$, $xy$             \\
          $A_1''$  & $+1$        & $+1$      & $+1$    & $-1$                &
          $-1$     & $-1$        & ---       & ---                           \\
          $A_2''$  & $+1$        & $+1$      & $-1$    & $-1$                &
          $-1$     & $+1$        & $z$       & ---                           \\
          $E''$    & $+2$        & $-1$      & $0$     & $-2$                &
          $+1$     & $0$         &$(R_x,R_y)$& $(xz, yz)$                    \\
          \bottomrule
        \end{tabular}
      \end{center}
    \end{remark}
    \item \ce{BF3} has several vibrational modes. Please prove step by step
    why the symmetric stretching mode is $A_1'$ mode. (\textbf{2 points})
    \item For all vibrational modes for \ce{BF3}, point out their IR and
    Raman activity. (\textbf{0.5 point})
  \end{enumext}
\end{problem}
\pagebreak
\begin{solution}\leavevmode
  \begin{enumext}
    \item
    \begin{enumext}
      \item Translational modes: $\Gamma_\text{trans} = A_2'' \oplus E'$.
      \item Rotational Modes: $\Gamma_\text{rot} = A_2' \oplus E''$.
      \item Vibrational Modes: Calculate $\Gamma_\text{tot}$ first.
      List the traces of the operators and \# of fixed atoms after operating.
      Then, the coefficients can be obtained via the inverse of the matrix
      \[
        \begingroup
        \renewcommand *\arraystretch{.8}
        \begin{tabular}{*4l}
          \toprule
          Op.         & Trace & \# Atoms & Product\\
          \midrule
          $1E$        & $3$   & $4$      & $12$   \\
          $2C_3$      & $0$   & $1$      & $0$    \\
          $3C_2'$     & $-1$  & $2$      & $-2$   \\
          $1\sigma_h$ & $1$   & $4$      & $4$    \\
          $2S_3$      & $-2$  & $1$      & $-2$   \\
          $3\sigma_v$ & $1$   & $2$      & $2$    \\
          \bottomrule
        \end{tabular}
        \endgroup\quad
        \bm\alpha =
          \begin{bmatrix}
            +1 & +1 & +2 & +1 & +1 & +2\\
            +1 & +1 & -1 & +1 & +1 & -1\\
            +1 & -1 &  0 & +1 & -1 &  0\\
            +1 & +1 & +2 & -1 & -1 & -2\\
            +1 & +1 & -1 & -1 & -1 & +1\\
            +1 & -1 &  0 & -1 & +1 &  0
          \end{bmatrix}^{-1}
          \begin{bmatrix}
            12\\0\\-2\\4\\-2\\2
          \end{bmatrix} =
          \begin{bmatrix}
            1\\1\\3\\0\\2\\1
          \end{bmatrix}.
      \]
      Then, $\Gamma_\text{tot} = \{1, 1, 3, 0, 2, 1\}$.
      The irreps. corresponding to vibrational mode is
      \[
        \Gamma_\text{vib}
      = \Gamma_\text{tot} - \Gamma_\text{trans} - \Gamma_\text{rot}
      = \{1, 0, 2, 0, 1, 0\} = A_1' \oplus 2E' \oplus A_2''.
      \]
    \end{enumext}
    \item \begin{proof}
      \# of bonds remain unmoved under the six operations are
      \begin{center}
        \begin{tabular}{*7l}
          \toprule
          Operation & $1E$        & $2C_3$ & $3C_2'$
                    & $1\sigma_h$ & $2S_3$ & $3\sigma_v$\\
          \midrule
          $N_\text{unmoved}$ & $3$ & $0$ & $1$ & $3$ & $0$ & $1$\\
          \bottomrule
        \end{tabular}
      \end{center}
      Substituting into the reduction formula
      $n_i = \frac1h \sum n_R \chi(R) \chi_i(R)$
      \begin{align*}
        n(A_1') & = \frac1{12}[(3)(1)(1) + (0)(2)(+1) + (1)(3)(1)
                             + (3)(1)(1) + (0)(2)(+1) + (1)(3)(1)] = 1,\\
        n(E')   & = \frac1{12}[(3)(1)(2) + (0)(2)(-1) + (1)(3)(0)
                             + (3)(1)(2) + (0)(2)(-1) + (1)(3)(0)] = 1,
      \end{align*}
      while other irreps. yield $0$.
      Therefore, $\Gamma_\text{stretch} = A_1' \oplus E'$.
    \end{proof}
    \item Due to the requirements of IR active and Raman active
    \begin{enumext}
      \item IR Active: The mode must transform as the linear functions
      ($x$, $y$, or $z$), as this indicates a change in the molecular
      dipole moment.
      \item Raman Active: The mode must transform as the quadratic functions
      ($x^2$, $y^2$, $z^2$, $xy$, $xz$, $yz$), as this indicates a change in
      the molecular polarizability.
    \end{enumext}
    Evaluate the three modes
    from $\Gamma_\text{vib} = A_1' \oplus 2E' \oplus A_2''$ by comparing with
    $D_{3h}$ group character table
    \begin{enumext}[columns = 2]
      \item $A_1'$: IR inactive, Raman active.
      \item $A_2''$: IR active, Raman inactive.
      \item $E'$: IR active, Raman active.
    \end{enumext}
  \end{enumext}
\end{solution}

\newpage

\begin{problem}
  \ce{W(CO)6} ($O_h$ point group) has three degenerate IR
  active vibrational modes with irreducible representation $T_{1u}$,
  as shown in the partial character table below
  \begin{center}
    \footnotesize
    \begin{tabular}{*{13}l}
      \toprule
      $O_h$            & $E$ & $8C_3$ & $6C_2$ & $6C_4$            &
      $3C_2 = (C_4)^2$ & $i$ & $6S_4$ & $8S_6$ & $3\sigma_h$       &
      $6\sigma_d$      & \parbox{7em}{Linear functions, rotations} &
      \parbox{4em}{Quadratic functions}\\
      \midrule
      $T_{1u}$         & $+3$& $0$    & $-1$   & $+1$ &
      $-1$             & $-3$& $-1$   & $0$    & $+1$ &
      $+1$             & $(x, y, z)$  & ---           \\
      \bottomrule
    \end{tabular}
  \end{center}
  Please follow the same following procedure and prove the character of $3C_2$
  ($3$ $C_2$-axes along $x$, $y$, and $z$ axes) is $-1$ for $T_{1u}$ irreps.
  (\textbf{1 point})
  \begin{snugshade}
    \paragraph{Procedure: $C_3$ about a body diagonal (111)}\leavevmode\par
    A representative $\qty{120}\degree$ rotation gives
    \[
      x \to y, \quad y \to z, \quad z \to x.
    \]
    So
    \[
      T_x \to T_y, \quad T_y \to T_z, \quad T_z \to T_x.
    \]
    Matrix:
    \[
      D^{(T_{1u})}(C_3) =
      \begin{pmatrix}
        0 & 0 & 1\\
        1 & 0 & 0\\
        0 & 1 & 0
      \end{pmatrix}.
    \]
    Character:
    \[
      \chi(C_3) = 0.
    \]
    So the $8C_3$ class has character 0.
  \end{snugshade}
  \begin{remark}
    The transformation of $T_x$, $T_y$, $T_z$ after $C_2(z)$
    operation, the corresponding matrix and the trace (the final value of
    character) are required.
  \end{remark}
\end{problem}
\begin{solution}
  A $C_2(z)$ operation represents a $\qty{180}\degree$ rotation about
  the $z$-axis, which gives
  \[
    x \to -x, \quad y \to -y, \quad z \to z.
  \]
  Consequently, the translational basis functions $(T_x, T_y, T_z)$
  transform as:
  \[
    T_x \to -T_x, \quad T_y \to -T_y, \quad T_z \to T_z.
  \]
  The transformation can be represented by the matrix
  \[
    D^{(T_{1u})}(C_2) =
    \begin{pmatrix}
    -1 & 0  & 0\\
    0  & -1 & 0\\
    0  & 0  & 1
    \end{pmatrix}.
  \]
  The character $\chi$ is defined as the trace of the
  representative matrix
  \[
    \chi(C_2) = \Tr(D^{(T_{1u})}(C_2)) =  (-1) + (-1) + 1 = -1.
  \]
  So, the $3C_2$ class in the $O_h$ point group has a character of $-1$ for
  the $T_{1u}$ irreducible representation.
\end{solution}
